\section{Related work, and what is claimed as new}\label{sec:related}

The result combines familiar ingredients, so its scope matters. The full-cone identity and the elementary quantifier inclusion are not claimed as new. The contribution is a finite-menu transition-certification construction, with explicit assumptions on the tube and scalarization, and a bounded account of how its verdict changes under selected control and disturbance extensions.

\subsection{Scalarization and finite-menu geometry}

Weighted-sum scalarization recovers supported outcomes of multi-objective problems and can miss unsupported points on a nonconvex frontier; alternative scalarizations address that limitation \citep{Koski1985,AthanPapalambros1996,DasDennis1997,Miettinen1999,Ehrgott2005,MessacMattson2002}. Theorem~\ref{thm:menu} has related convex geometry, but its exact statement is a margin-first acceptance result: the required-margin vector is formed first and the scalarized requirement is assumed to be $w\cdot d_a$. It compares the up-closed convex hull of a finite menu with the union of the menu's orthants. This is not a new scalarization theorem, nor does it establish that every exact-tube scalarization has that form; Lemma~\ref{lem:sync} states when the two path orders agree.

The quantifier gap itself can be static. For example, the two alternative margin vectors $(-4/5,6/5)$ and $(6/5,-4/5)$ are not componentwise nonnegative, yet every nonzero nonnegative weight has at least one nonnegative alternative. In this paper, paths and disturbances specify how such margins arise and add transition-specific questions about transient breaches, terminal conditions, and implementable controls. They are an application setting, not a premise that makes the quantifier mechanism inherently dynamic. The relevant distinction from Pareto-front analysis is the decision object: acceptance of a state under a menu and a rule, rather than recovery of a particular outcome by a weighted-sum optimizer. The convex geometry is shared; the questions are not identical.

\subsection{Adjustable robustness, randomization, and physical blending}

The order $\exists a\,\forall w$ versus $\forall w\,\exists a_w$ resembles the here-and-now/wait-and-see distinction in adjustable robust optimization \citep{BenTal2004,BenTal2009}. The analogy is limited: here $w$ is an assessment weighting, not a realized uncertainty, and the result is a feasibility comparison, not a theorem that reverses the usual ordering of static and adjustable policy classes. On a common model and admissible-policy class, allowing adaptation includes the static policies; comparing a fractional physical control with a compulsory sequence of complete excursions compares different control extensions.

Fractional blending in Theorem~\ref{thm:blend} changes the physical control by allocating the primitive flows fractionally; it is not a lottery over primitive plans. In the shared-shock example at the canonical point, the best primitive worst margin is $-4/5$, while the balanced blend has worst margin $-1/20$: blending strictly improves this maximin margin but does not restore feasibility. We therefore describe the result as \emph{failure of feasibility restoration} at that point, not as randomization-proofness. Randomized decisions and their worst-case values are studied separately in robust optimization \citep{Delage2019,KrauseRoperGolovin2011,VorobeychikLi2014,Sinha2018,Sessa2020}.

Nor is the blend-versus-sequencing comparison a counterexample to Kuhn's realization-equivalence theorem. Under its hypotheses, mixed and behavioral strategies induce the same distribution over paths; equal path laws preserve probabilities of measurable safety events and distributional worst-case criteria. A deterministic fractional control instead produces a different physical path. The complete-excursion proposition below is conditional on its specified union-of-tubes control class. Under a distinct natural derivative implementation, a three-switch schedule can keep both canonical margins positive; this is not a counterexample to that conditional proposition, but it rules out a broad claim that all finite-rate switching fails. Any literal mixed-strategy interpretation would require a separate payoff, information structure, and disturbance model. Main and Randour study finite-memory strategy restrictions, but this manuscript does not establish that its control problem is an instance of their automaton model \citep{Kuhn1953,MainRandour2024}.

\subsection{Weak and strong sustainability}

The distinction formalized here is related to the weak-sustainability treatment of substitution among forms of capital and the strong-sustainability treatment of selected critical stocks as separately binding \citep{PearceAtkinson1993,Beckerman1994,Daly1995,Ayres2001,Neumayer2013,DietzNeumayer2007}. Both literatures include formal and empirical work; we do not characterize the wider debate as merely verbal, and we do not rank the doctrines. The operators idealize one transition-certification rule under each reading. The result establishes that these particular rules can disagree on a declared menu and disturbance class, not that any boundary or criterion ought normatively to be noncompensatory.

\subsection{Bounded contribution}

The paper's contribution is a source-checkable construction with four layers. First, the full-cone identity locates the gap in plan selection rather than loss of information at a fixed trajectory. Second, the finite-menu theorem gives the convex-margin description under its explicit margin-first hypothesis, and a rational witness realizes a nonempty transition-certification gap with an exact reserve threshold. Third, the paper analyzes a closed fractional-flow extension and several separately declared disturbance classes; these results are conditional on their respective control and shock models, not on a general equivalence between blending and randomization. Fourth, the CES ladder gives exact diagonal comparisons under the stated collapse convention and normalization; it does not identify a calibrated fishery index. The Northern-cod comparison is a historical dimensional anchor only.

We therefore claim neither priority for the quantifier-order observation, the standard convex-hull separation, nor the general usefulness of randomization. The claim is the explicit, reproducible scope of this particular transition-assessment construction and the boundaries of the conclusions that follow from it.
